%==============================================================================
% Sjabloon onderzoeksvoorstel bachproef
%==============================================================================
% Gebaseerd op document class `hogent-article'
% zie <https://github.com/HoGentTIN/latex-hogent-article>

% Voor een voorstel in het Engels: voeg de documentclass-optie [english] toe.
% Let op: kan enkel na toestemming van de bachelorproefcoördinator!
\documentclass{hogent-article}

% Invoegen bibliografiebestand
\addbibresource{voorstel.bib}

% Informatie over de opleiding, het vak en soort opdracht
\studyprogramme{Professionele bachelor toegepaste informatica}
\course{Bachelorproef}
\assignmenttype{Onderzoeksvoorstel}
% Voor een voorstel in het Engels, haal de volgende 3 regels uit commentaar
% \studyprogramme{Bachelor of applied information technology}
% \course{Bachelor thesis}
% \assignmenttype{Research proposal}

\academicyear{2025-2026} % TODO: pas het academiejaar aan

% TODO: Werktitel
\title{Van VPN naar Zero-Trust: Reductie van laterale beweging in IT-netwerken van woonzorgcentra}

% TODO: Studentnaam en emailadres invullen
\author{Egon Wiedler}
\email{Egon.Wiedler@student.hogent.be}

% TODO: Medestudent
% Gaat het om een bachelorproef in samenwerking met een student in een andere
% opleiding? Geef dan de naam en emailadres hier
% \author{Yasmine Alaoui (naam opleiding)}
% \email{yasmine.alaoui@student.hogent.be}

% TODO: Geef de co-promotor op
\supervisor[Co-promotor]{S. Beekman (Synalco, \href{mailto:sigrid.beekman@synalco.be}{sigrid.beekman@synalco.be})}

% Binnen welke specialisatierichting uit 3TI situeert dit onderzoek zich?
% Kies uit deze lijst:
%
% - Mobile \& Enterprise development
% - AI \& Data Engineering
% - Functional \& Business Analysis
% - System \& Network Administrator
% - Mainframe Expert
% - Als het onderzoek niet past binnen een van deze domeinen specifieer je deze
%   zelf
%
\specialisation{System \& Network Administrator}
\keywords{Zero Trust Network Access (ZTNA), Laterale beweging, Micro-segmentatie, IoMT, Open-source, VPN }

\usepackage{pgfgantt}
\begin{document}

\begin{abstract}
	
De toenemende digitalisering binnen de zorgsector, gecombineerd met de noodzaak voor externe toegang tot patiëntgegevens en medische IoT-apparatuur, creëert een complex beveiligingsvraagstuk voor hedendaagse woonzorgcentra. een speciefiek en zeer kritisch probleem hierbinnen is de architecturale zwakte van traditionel VPN-oplossingen. Deze verlenen na authenticatie vaak onnodig brede netwerktoegang en zullen zo de laterale beweging van kwaadwillende naar cruciale systemen faciliteren. Dit onderzoeksvoorstel zal zich daarom centreren rond de hoofdonderzoeksvraag: “In welke mate reduceert de implementatie van een open-source Zero-Trust Network Access (ZTNA) oplossing het aanvalsoppervlak voor laterale beweging binnen de IT-infrastructuur van een woonzorgcentrum, vergeleken met een traditionele VPN?”. De voorgestelde methodologie omvat het opzetten van een representatieve gesimuleerde omgeving ook wel een Proof of Concept, gebasseerd op een segment van bestaande netwerkinfrastructuur van een zorginstelling. waarin zowel de traditionele VPN methode en de ZTNA-oplossing technisch worden geïmplementeerd. Vervolgens zal er door kwalitatieve netwerkanalyses en bereikbaarheidstesten, het verschil op open netwerkpaden en het aanvalsoppervlak tussen beide scenario’s gemeten worden. De verwachte resultaten duiden zich op een significante en meetbare reductie van de interne bereikbaarheid in het ZTNA-scenario zonder verlies van noodzakelijke functionalitiet. Als de resultaten doormiddel van het ondeerzoek bevestigd worden, zal kunnen geconcludeerd worden dat de implementatie van open-source ZTNA-technologie een superieure en noodzakelijke strategie is om de weerbaarheid van zorgnetwerken tegen interne verspreiding van cyberdreigingen.


\end{abstract}

\tableofcontents

% De hoofdtekst van het voorstel zit in een apart bestand, zodat het makkelijk
% kan opgenomen worden in de bijlagen van de bachelorproef zelf.
%---------- Inleiding ---------------------------------------------------------

% TODO: Is dit voorstel gebaseerd op een paper van Research Methods die je
% vorig jaar hebt ingediend? Heb je daarbij eventueel samengewerkt met een
% andere student?
% Zo ja, haal dan de tekst hieronder uit commentaar en pas aan.

%\paragraph{Opmerking}

% Dit voorstel is gebaseerd op het onderzoeksvoorstel dat werd geschreven in het
% kader van het vak Research Methods dat ik (vorig/dit) academiejaar heb
% uitgewerkt (met medesturent VOORNAAM NAAM als mede-auteur).
% 


\section{Inleiding}%
\label{sec:inleiding}

\subsection{Situering}
IT-infrastructuur binnen de zorgsector ondergaat een snelle transformatie. Waar een woonzorgcentrum (WZC) vroeger volstond met enkele administratieve werkplekken, is er vandaag sprake van een complex ecosysteem met elektronische patiëntendossiers (EPD), mobiele zorgtoepassingen en een groeiend aantal medische IoT-apparaten (``Internet of Medical Things''). Deze digitalisering brengt echter nieuwe beveiligingsrisico's met zich mee, zeker gezien de noodzaak voor externe toegang door leveranciers, artsen en thuiswerkend personeel.

Dit onderzoek vertrekt vanuit een concrete casus bij woonzorgcentrum Curando. Binnen deze organisatie wordt momenteel gebruikgemaakt van een traditionele VPN-oplossing (Virtual Private Network) om externe toegang te faciliteren. Hoewel deze technologie de verbinding versleutelt, hanteert ze in de praktijk vaak een perimeter-gebaseerd beveiligingsmodel: eenmaal een gebruiker geauthenticeerd is, krijgt deze toegang tot een groot deel van het interne netwerksegment.

\subsection{Probleemstelling}
Het specifieke probleem dat zich bij deze casus stelt, is het gebrek aan interne segmentatie na authenticatie. Een gecompromitteerd toestel van een externe medewerker kan in de huidige netwerkarchitectuur als springplank dienen om andere servers en kritieke systemen te bereiken. Dit fenomeen, bekend als \emph{laterale beweging}, vormt een aanzienlijk risico voor de bedrijfscontinuïteit en patiëntgegevens, zeker in het licht van recente ransomware-aanvallen op zorginstellingen. De huidige VPN-architectuur biedt onvoldoende granulariteit om dit risico effectief in te perken zonder de operationele werking te verstoren. Er is daarom nood aan een moderne benadering gebaseerd op Zero-Trust principes (``never trust, always verify'').

\subsection{Onderzoeksvraag}
Om een oplossing te bieden voor dit veiligheidsrisico, staat in dit onderzoek de volgende hoofdvraag centraal:

\begin{quote}
	\textbf{``In welke mate reduceert de implementatie van een open-source Zero-Trust Network Access (ZTNA) oplossing het aanvalsoppervlak voor laterale beweging binnen de IT-infrastructuur van een woonzorgcentrum, vergeleken met een traditionele VPN?''}
\end{quote}

Ter beantwoording van deze hoofdvraag wordt het onderzoek opgesplitst in twee domeinen, die elk worden uitgewerkt aan de hand van specifieke deelvragen:

\begin{itemize}
	\item \textbf{Probleemdomein:} Dit domein richt zich op het analyseren van de huidige netwerkarchitectuur en de bijbehorende beveiligingsrisico’s binnen zorgomgevingen. De volgende deelvragen staan hierbij centraal:
	\begin{enumerate}
		\item Wat zijn de fundamentele architecturale verschillen tussen VPN en Zero Trust Network Access (ZTNA) volgens het NIST-raamwerk?
		\item Welke specifieke beveiligingsrisico’s ontstaan bij het gebruik van traditionele VPN-oplossingen in een zorgnetwerk met IoT-apparatuur?
		\item Hoe groot is het aanvalsoppervlak van een traditioneel VPN-gebaseerd zorgnetwerk en welke kwetsbaarheden zijn hierbij het meest relevant?
	\end{enumerate}
	
	\item \textbf{Oplossingsdomein:} Dit domein focust op het evalueren van alternatieve netwerktoegangsmodellen en het praktisch aantonen van hun meerwaarde. Dit leidt tot de volgende deelvragen:
	\begin{enumerate}
		\item Welke open-source ZTNA-oplossing (zoals NetBird of Headscale) is het meest geschikt voor toepassing in een KMO-omgeving?
		\item Op welke manier kan een representatieve Proof of Concept (PoC) worden opgezet om VPN- en ZTNA-architecturen objectief met elkaar te vergelijken?
		\item Hoe kan het verschil in aanvalsoppervlak tussen een VPN- en een ZTNA-oplossing binnen deze PoC kwantitatief meetbaar worden gemaakt?
	\end{enumerate}
\end{itemize}


\subsection{Onderzoeksdoelstelling}
Het doel van deze bachelorproef is niet enkel theoretisch; het beoogt tastbare, technisch meetbare resultaten. Het concrete eindresultaat is een Proof of Concept (PoC) in een gesimuleerde omgeving, die de netwerkstructuur van de betreffende zorginstelling nabouwt. Binnen deze omgeving worden de traditionele VPN en een open-source ZTNA-oplossing (zoals Headscale of NetBird) technisch geïmplementeerd en vergeleken.

Het succes van dit onderzoek wordt bepaald door een kwantitatieve analyse waarin wordt aangetoond of, en in welke mate, de ZTNA-oplossing het aantal bereikbare netwerkpaden reduceert ten opzichte van de nulmeting (VPN). Dit resulteert uiteindelijk in een technisch adviesrapport voor de doelgroep, waarmee zij een gefundeerde beslissing kunnen nemen over een eventuele implementatie van deze architectuur.
%---------- Stand van zaken ---------------------------------------------------

\section{Literatuurstudie}%
\label{sec:literatuurstudie}

Deze literatuurstudie onderzoekt de fundamentele principes van Zero Trust Network Access (ZTNA), de architecturale componenten, de noodzaak voor migratie vanaf traditionele VPN’s en de specifieke toepassing binnen de gezondheidszorg en open-source oplossingen.

\subsection{Definitie en Kernprincipes}
Zero Trust (ZT) wordt niet gedefinieerd als een enkele technologie, maar als een cybersecurity-paradigma gericht op de bescherming van resources. Het kernprincipe, zoals gedefinieerd door het National Institute of Standards and Technology (NIST), is dat vertrouwen nooit impliciet wordt verleend op basis van netwerklocatie of het eigendom van een apparaat \parencite{rose2020zero}. Een Zero Trust Architectuur (ZTA) vereist dat elke toegang tot bedrijfsbronnen per sessie wordt geauthenticeerd en geautoriseerd, waarbij gebruik wordt gemaakt van dynamisch beleid om onzekerheid bij toegangsbeslissingen te minimaliseren \parencite{rose2020zero}.

\subsection{Architectuur: Control Plane vs. Data Plane}
Een essentieel kenmerk van moderne ZTNA-platformen is de strikte scheiding tussen de \textit{control plane} (besturingslaag) en de \textit{data plane} (gegevenslaag).
\begin{itemize}
	\item \textbf{Control Plane:} werkt als het "brein" van het systeem. Het verifieert de identiteit, controleert de status van het apparaat (posture) en evalueert het beleid voordat er data wordt uitgewisseld \parencite{terrazone2025}.
	\item \textbf{Data Plane:} werkt als het "lichaam" dat de daadwerkelijke data transporteert via versleutelde tunnels, maar pas nadat de control plane toestemming heeft gegeven \parencite{terrazone2025}.
\end{itemize}
Volgens de NIST-standaard bestaat de logische architectuur uit een \textit{Policy Engine} (PE) die beslissingen neemt, een \textit{Policy Administrator} (PA) die de verbinding opzet, en een \textit{Policy Enforcement Point} (PEP) dat de verbinding daadwerkelijk toestaat of blokkeert \parencite{rose2020zero}. Deze scheiding zorgt voor schaalbaarheid en veiligheid; als de control plane uitvalt, kunnen bestaande sessies op de data plane vaak blijven bestaan \parencite{terrazone2025}.

\subsection{Verschuiving van VPN naar ZTNA}
Traditionele VPN-technologie, gebaseerd op het ``kasteel en gracht''-model, voldoet niet meer in de moderne hybride werkomgeving. VPN's verlenen vaak toegang tot het volledige netwerk na authenticatie, wat \textit{laterale verplaatsing} door aanvallers mogelijk maakt en het risico op ransomware vergroot \autocite{barracuda2024}. Daarnaast zorgen VPN's vaak voor prestatieproblemen door verkeer onnodig via een centraal punt te leiden ("hairpinning") \autocite{seqrite2025}.

ZTNA keert het model om van "verbinden dan verifiëren" naar "verifiëren dan verbinden" \autocite{nflo2025}. Gartner voorspelt dat in 2025 minstens 70\% van de nieuwe implementaties voor toegang op afstand gebruik zal maken van ZTNA in plaats van VPN, gedreven door de toename van hybride werken en cloud-adoptie \autocite{gartner2022}.

\subsection{Toepassing in Gezondheidszorg}
De zorgsector is een belangrijk doelwit voor cyberaanvallen. Zorginstellingen kampen met verouderde apparatuur, het \textit{Internet of Medical Things} (IoMT), en strenge compliance-eisen zoals HIPAA \parencite{anacyber2024}. Z-CERT meldt in het dreigingsbeeld van 2024 een toename van digitale aanvallen, waaronder ransomware en datalekken, waarbij clouddiensten steeds vaker het doelwit zijn \parencite{zcert2024}.

ZTNA biedt hier een oplossing door netwerksegmentatie toe te passen en apparaten onzichtbaar te maken voor het publieke internet \parencite{Wal2022}. Het faciliteert veilige overdracht van kritieke data, zoals radiologiebeelden, zonder de workflow te vertragen en dwingt \textit{least-privilege} toegang af \parencite{appgate2023}. Dit is cruciaal om patiëntendossiers te beschermen tegen gijzelsoftware en datadiefstal \parencite{hoop2024}.

\subsection{Implementatie}
De markt biedt diverse open-source tools voor ZTNA-implementaties:
\begin{itemize}
	\item \textbf{OpenZiti:} Richt zich op het inbedden van Zero Trust direct in applicaties via SDK's, waardoor een overlay-netwerk ontstaat dat geen open poorten vereist \parencite{aimultiple2025}. Het verbindt services in plaats van machines \parencite{OpenZiti2025}.
	\item \textbf{Pomerium:} Maakt gebruik van een \textit{identity-aware proxy} architectuur die elke actie continu verifieert \autocite{aimultiple2025}.
	\item \textbf{Octelium:} Een recenter platform dat functioneert als een "unified zero trust secure access platform". Het combineert functionaliteiten van een remote access VPN, ZTNA, en API-gateway, en ondersteunt protocollen zoals WireGuard en QUIC \autocite{octelium2025}.
	\item \textbf{Pritunl Zero \& OpenVPN:} Bieden eveneens ZTNA-functionaliteiten, met focus op respectievelijk SSH-management en netwerktunneling \autocite{aimultiple2025}.
\end{itemize}

%---------- Methodologie ------------------------------------------------------
\section{Methodologie}%
\label{sec:methodologie}

Dit onderzoek wordt gekenmerkt als een \textbf{vergelijkende} studie ondersteund door een technische \textbf{Proof of Concept (PoC)}. Het doel is om op een objectieve manier te bekijken hoe de verschillen zijn van een netwerk met ZTNA-principes en een netwerk die gebruikt maakt van een traditionele VPN.

Om de hoofdvraag te beantwoorden, wordt het onderzoek opgesplitst in vier fasen. In elk van deze fases zal een specifieke onderzoekstechniek worden gehanteerd.

\subsection{Fase 1: Literatuurstudie}
In de eerste fase wordt middels een \textbf{literatuurstudie} de theoretische basis gelegd. Hierbij wordt gebruikgemaakt van het NIST SP 800-207 raamwerk \parencite{rose2020zero} om de architecturale vereisten van Zero Trust vast te leggen en meer inzicht te geven.
Daarnaast wordt een \textbf{vergelijkende analyse} uitgevoerd van beschikbare open-source ZTNA-oplossingen (zoals Headscale, NetBird of Firezone). De selectie van de tool voor de PoC gebeurt op basis van vooraf gedefinieerde criteria zoals bv.:
\begin{itemize}
	\item Compatibiliteit met het WireGuard-protocol.
	\item Granulariteit van Access Control Lists (ACL's).
	\item Geschiktheid voor implementatie in een KMO-omgeving (zoals een woonzorgcentrum).
\end{itemize}

\subsection{Fase 2: Proof of Concept}
De kern van de methodologie is de \textbf{simulatie} van de IT-infrastructuur. Er wordt een testomgeving opgezet die de netwerktopologie van de casus (het woonzorgcentrum) nabouwt.
\begin{enumerate}
	\item \textbf{Architectuur:} De omgeving wordt gevirtualiseerd en gesegmenteerd in verschillende zones: een administratief segment, een medisch segment (EPD/IoT) en een gastennetwerk.
	\item \textbf{Scenario A (Nulmeting):} Implementatie van een traditionele VPN-server (bijv. OpenVPN of WireGuard in gateway-modus) die toegang geeft tot het netwerk, representatief voor de huidige situatie.
	\item \textbf{Scenario B (Experiment):} Implementatie van de geselecteerde ZTNA-oplossing, waarbij /\newline icro-segmentatie wordt toegepast op basis van identiteit in plaats van netwerklocatie.
\end{enumerate}


\subsection{Fase 3: Experiment en Meting}
Om de onderzoeksvraag meetbaar te maken, wordt een \textbf{experiment} uitgevoerd volgens het ``Assume Breach''-principe. Hierbij wordt aangenomen dat het endpoint van een externe medewerker is gecompromitteerd en beschikt over geldige toegangsrechten.

\begin{itemize}
	\item \textbf{Netwerkexposure (nulmeting):} Vanuit het gecompromitteerde endpoint worden geautomatiseerde netwerkscans uitgevoerd met tools zoals \textbf{Nmap} om het aantal bereikbare hosts en openstaande servicepoorten (TCP/UDP) in andere netwerksegmenten te bepalen.
	
	\item \textbf{Laterale beweging:} Er wordt getest in welke mate het mogelijk is om vanuit het gecompromitteerde endpoint verbinding te maken met niet-geautoriseerde systemen en diensten (zoals SMB-, SSH- of HTTP-services) binnen het interne netwerk.
	
	\item \textbf{Identiteits- en contextafhankelijkheid:} De experimenten worden herhaald met identieke gebruikerscredentials maar vanuit een ander device of netwerkcontext, om het verschil in toegangscontrole tussen VPN en ZTNA te analyseren.
	
	\item \textbf{Blast radius:} De effectieve toegang tot netwerkresources wordt vergeleken met de vooraf gedefinieerde toegangsrechten, om de impact van een succesvolle initiële compromittering te evalueren.
	
	\item \textbf{Validatie:} Met \textbf{Wireshark} en logbestanden van de gebruikte oplossingen wordt gecontroleerd of het verkeer via de correcte tunnels verloopt en of toegangsregels correct worden afgedwongen.
\end{itemize}


\subsection{Fase 4: Vergelijkende Analyse}
De data uit Fase 3 (de scanresultaten van Scenario A versus Scenario B) worden naast elkaar gelegd. De \textbf{kwantitatieve reductie} van het aanvalsoppervlak wordt berekend. Op basis hiervan wordt een advies geformuleerd over de haalbaarheid en de veiligheidswinst voor het woonzorgcentrum.


\subsection{Tijdsplanning en Deliverables}
De uitvoering van dit onderzoek is gepland over een periode van 14 weken.

\begin{figure}[ht]
	\centering
	% Resizebox dwingt de grafiek om exact in de kolom te passen
	\resizebox{\columnwidth}{!}{%
		\begin{ganttchart}[
			vgrid,
			hgrid,
			x unit=0.45cm,       % Smalle eenheden zodat het past
			y unit chart=0.5cm, 
			title label font=\bfseries\tiny, % Heel klein lettertype voor de cijfers
			bar label font=\tiny,            % Heel klein lettertype voor de taken
			bar/.append style={fill=black!30},
			bar height=0.6,
			group right shift=0,
			group top shift=0.7,
			group height=.3
			]{1}{14}
			
			% We korten de titel in om ruimte te besparen
			\gantttitle{Weken (Feb-Mei)}{14} \\
			\gantttitlelist{1,...,14}{1} \\
			
			% De taken
			\ganttbar{1. Lit. \& Selectie}{1}{3} \\
			\ganttbar{2. Lab \& VPN}{4}{6} \\
			\ganttbar{3. Impl. ZTNA}{7}{9} \\
			\ganttbar{4. Testen}{5}{12} \\
			\ganttbar{5. Analyse}{13}{14} \\
			
			\ganttbar[bar/.append style={fill=black!60}]{Documentatie}{1}{14}
			
		\end{ganttchart}%
	}
	\caption{Projectplanning}
	\label{fig:planning}
\end{figure}

%---------- Verwachte resultaten ----------------------------------------------
\section{Verwacht resultaat, conclusie}%
\label{sec:verwachte_resultaten}


\subsection{Verwacht resultaat}



De primaire verwachting van dit onderzoek is dat de implementatie van Zero-Trust Network Access (ZTNA) het aantal bereikbare netwerkpaden drastisch zal reduceren ten opzichte van de traditionele VPN-configuratie. Waar bij de VPN-nulmeting wordt verwacht dat een gecompromitteerd toestel onbeperkte of minimale beperkte toegang heeft tot achterliggende servers en IoT-apparatuur, zal de ZTNA-oplossing deze zichtbaarheid technisch beperken tot enkel de strikt noodzakelijke applicatieservices. Dit verschil in aanvalsoppervlak zal kwantitatief worden aangetoond middels netwerkscans en gevisualiseerd worden in een vergelijkende grafieken en tabellen.

\subsection{Conclusie}
Indien de meetresultaten de hypothese bevestigen, luidt de conclusie dat open-source micro-segmentatie een noodzakelijke en effectieve verdedigingslinie vormt tegen laterale beweging binnen zorgnetwerken. Dit biedt IT-beheerders een concreet bewijs dat het isoleren van apparaten via ZTNA de impact van ransomware-infecties beperkt tot het individuele toestel, waardoor de bedrijfscontinuïteit en patiëntveiligheid behouden blijven zonder dat hierbij hoge licentiekosten voor enterprise-software noodzakelijk zijn.


\printbibliography[heading=bibintoc]

\end{document}